\section{Introduction}


% 随着自动驾驶的普及，随着无人驾驶车辆越来越多，无人车和有人车混在一起。这个混合交通场景正在成为未来的发展趋势。

%混合驾驶动态路由问题作为这个场景的核心问题，正在越来越受到关注。

% 问题的定义是给定OD，动态地选择每一辆自动驾驶汽车的路径（控制），在给定人类驾驶汽车行驶背景的情况下（环境），最小化全局通行时间（优化目标）。

% 在场景中，能够调度的是CAV，HDV的行为间接地受CAV决策的控制。CAV和HDV相互耦合，我们可以直接控制CAV，通过CAV产生的拥堵，间接地影响HDV。HDV 的影响又会改变决策环境，需要根据环境变化调整CAV的决策。

% 传统建模方法把微观的路径选择，汇集成宏观的交通流，然后通过交通流理论分配各个OD的路由，他的优化目标是使得城市的交通流达到交通均衡。通过求解均衡太优化个体的决策。问题主要是忽略了微观信息，或者hand-craft（建模的过程会有信息损失）。通过宏观交通分配的原则指导个体的交通决策。

%传统建模方法建模HDV的交通流，构造背景环境


% 这两种方法的问题，需要对HDV行为进行显式建模，利用规则对HDV进行建模。（一定要围绕这种混合路由的问题来讲）这两种方法HDV都是ruler driven的。

% 第三个问题是强化学习计算复杂度很高，需要昂贵的环境交互。


% To bridge this gap，我们提出了基于世界模型的方法，他的好处是什么，现在已经在哪些领域取得成功。世界模型强化学习可以解决上述问题，通过data-driven的方式，建立这种HDV+CAV的混合环境转移的预测模型。然后通过强化学习训练优化决策。好处1:rule-based 到 data-driven。 好处2:不光仿真hdv到行为，还会仿真hdv和cav交互关系，使其得到很好的建模下（建模hdv-cav的混合环境）。好处3:预测agent行为下交通环境的变化，减少训练的交互成本和仿真的开销。

Dynamic routing on urban road networks is a sequential decision problem in which connected and autonomous vehicles (CAVs) and human-driven vehicles (HDVs) repeatedly choose routes toward their destinations, and every choice reshapes the congestion faced by all later decisions. This coupling is asymmetric, because CAVs can be centrally coordinated whereas self-interested HDVs cannot be rerouted directly and instead adapt to the congestion patterns created by the coordinated vehicles. Coordinating this limited controllable fleet to reduce the mean travel time of both populations is therefore a difficult dynamic routing game \cite{calderone2017markov,lazar2020routing}, whose solution is increasingly critical for urban traffic management as CAV adoption continues to grow.

In the literature, research on this problem has followed three main directions, namely analytical formulations, model-free multi-agent reinforcement learning (MARL), and world-model-based learning. Analytical models \cite{tanaka2020linearly,cabannes2022solving,kolarich2022stackelberg} offer interpretable descriptions of traffic externalities and HDV response, but depend on prescribed link-cost functions, behavioral assumptions, or repeated equilibrium computation, which become restrictive once time-varying departures, incidents, and microscopic vehicle interactions jointly determine the traffic transition.
Model-free MARL avoids these assumptions by letting CAVs learn state-dependent route policies in microscopic simulators, formulating the problem as a Markov game with executable urban environments \cite{lazar2021learning,shou2022multi,akman2025routerl,akman2026urb}. This flexibility, however, is bought with repeated environment queries, and the delayed effect of route choices on shared bottlenecks makes real interaction expensive to reuse for long-horizon, multi-vehicle credit assignment.
World-model reinforcement learning removes this cost by learning latent dynamics from collected trajectories and updating policies in imagined rollouts \cite{hafnerdream,hafnermastering,hafner2025mastering}, a paradigm already validated in driving and multi-agent domains \cite{mamba,li2024think2drive,yang2026raw2drive}. Porting it to mixed-autonomy dynamic routing, however, is not straightforward, because the problem couples a graph-structured spatiotemporal state with vehicle-local route semantics, an uncontrolled HDV background whose demand varies across episodes, and per-vehicle actions tied to a network-wide objective.

In summary, analytical models are interpretable but inflexible, model-free MARL is flexible but sample-inefficient, and existing world models, though sample-efficient, are not built for mixed-autonomy routing. A world model that jointly captures graph-level traffic transitions, HDV demand, and vehicle-to-network credit assignment is still missing, and building one is precisely the goal of this work.

To this end, we propose UrbanDreamer, a world-model-based MARL framework designed around the three interfaces above. A parameter-shared actor retains the executable control semantics of the routing problem, in which each active CAV selects a candidate route from its current segment to its destination. UrbanDreamer then projects the sampled route choices onto a planned CAV pressure field defined over road-segment nodes, placing vehicle actions and network dynamics in a common graph coordinate system. On top of this interface, a graph convolutional encoder represents the traffic field, and a state-to-state flow-matching model predicts its next latent representation conditioned on CAV pressure, departures, and a history-conditioned estimate of HDV route demand. During imagined rollouts, a deterministic vehicle wrapper preserves vehicle identities, route progress, departures, and arrivals, while centralized per-vehicle and global critics supply individual and network-level value signals. Simulator collection episodes periodically refresh the world model, and held-out simulator episodes serve as the basis for policy evaluation. Our contributions are summarized as follows.
\begin{itemize}
\item We propose a world-model-based MARL framework for mixed-autonomy dynamic routing, in which additional CAV policy updates are performed in learned latent traffic dynamics, substantially reducing the reliance on repeated microscopic simulation.
\item We design a graph-aligned action interface and traffic world model that lift executable per-CAV route choices onto segment-level latent transitions, combining graph convolutional traffic encoding with conditioned state-to-state flow matching and a history-conditioned, slow-timescale model of aggregate HDV demand.
\item We develop identity-preserving imagined route learning that couples a parameter-shared CAV actor with centralized per-vehicle and global critics, preserving vehicle-specific credit under a network-wide objective through a two-stage, simulator-grounded pipeline alternating online collection, world-model adaptation, and imagined actor-critic optimization.
\end{itemize}
Experiments on three urban road networks with 636, 1,035, and 2,068 scheduled vehicles, instantiated through the RouteRL and URB framework, show that UrbanDreamer reaches strong routing policies with substantially fewer real simulator interactions and lower wall-clock training time than model-free MARL baselines, attaining the lowest global mean travel time on two of the three networks.
